\section*{Capítulo \ref{ch:b3:molecular-evolution} --- Evolução molecular e filogenômica}

\begin{solution}{exo:b3:molecular-evolution:1}
$k = \mu$: a taxa neutra de substituição por sítio por geração é igual à
taxa de mutação. A cada geração surgem $2N\mu$ mutações neutras novas e
cada uma tem probabilidade $1/2N$ de se fixar; uma população maior faz
mais mutações e dá a cada uma uma chance proporcionalmente menor, e os
dois efeitos se cancelam exatamente.
\end{solution}

\begin{solution}{exo:b3:molecular-evolution:2}
\omterm{def:b3:genomics:comparative}{Ortólogos}: genes de duas espécies descendentes de um gene do ancestral
comum --- a $\beta$-globina humana e a $\beta$-globina de camundongo.
\omterm{def:b3:genomics:comparative}{Parálogos}: genes de um mesmo genoma descendentes de uma duplicação --- a
$\alpha$- e a $\beta$-globina humanas, ou a $\beta$ e a $\gamma$.
\omterm{def:b3:molecular-evolution:duplication}{Pseudogene}: uma cópia decaída que já não codifica proteína --- o gene
$\psi\beta$ do agrupamento da $\beta$-globina humana, completo com
códons de parada.
\end{solution}

\begin{solution}{exo:b3:molecular-evolution:3}
$\omega$ compara a \omterm{thm:b3:molecular-evolution:neutral}{taxa de substituição} que muda aminoácidos com a \omterm{thm:b3:molecular-evolution:neutral}{taxa
de substituição} silenciosa, que é a linha de base neutra. $0.02$: \omterm{def:b3:molecular-evolution:dnds}{seleção
purificadora} forte, com \qty{98}{\%} das mudanças de aminoácido removidas
--- actina, \omterm{def:b3:chromatin-epigenetics:nucleosome}{histona} H3. $1.0$: nenhuma restrição --- um \omterm{def:b3:molecular-evolution:duplication}{pseudogene}.
$2.5$: \omterm{def:b3:molecular-evolution:dnds}{seleção positiva}, com as mudanças de aminoácido favorecidas --- a
fenda de ligação a peptídeos do MHC, o \omterm{def:b3:virology:virus}{envelope} do HIV, a lisozima de
ruminante.
\end{solution}

\begin{solution}{exo:b3:molecular-evolution:4}
As proteínas comparadas entre mamíferos mudam a cerca de uma substituição
por sítio por bilhão de anos; num genoma de um bilhão de sítios
codificadores e numa geração de alguns anos, isso é da ordem de uma
substituição fixada por geração, ou uma a cada dois anos. O custo da
seleção de Haldane permite cerca de uma substituição selecionada a cada
trezentas gerações. As taxas diferem por um fator de cem ou mais, de modo
que quase todas as mudanças fixadas precisam ser neutras.
\end{solution}

\begin{solution}{exo:b3:molecular-evolution:5}
Neutra: $1/2N = 10^{-4}$. $s = 0.001$, $4Ns = 20$: $(1 - \mathrm{e}^{-0.002})/
(1 - \mathrm{e}^{-20}) = 0.002$. $s = 0.01$:
$0.0198$. $s = -0.001$: $(1 -
\mathrm{e}^{0.002})/(1 - \mathrm{e}^{20}) = 0.002\,\mathrm{e}^{-20} =
4\times 10^{-12}$. Nenhuma é efetivamente
neutra: isso exigiria $|4Ns|
\lesssim 1$, isto é, $|s| \lesssim 5\times 10^{-5}$; um coeficiente de seleção de um décimo de por cento já é
decisivo numa população de cinco mil.
\end{solution}

\begin{solution}{exo:b3:molecular-evolution:6}
Sem correção: $T = 0.12/(2\times 2\times 10^{-9}) = \qty{30}{Myr}$.
Jukes--Cantor: $d = -\tfrac{3}{4}\ln(1 - 0.16) = 0.131$, $T =
\qty{33}{Myr}$. A correção chega a um fator de dois quando
$-\tfrac{3}{4}
\ln(1 - 4p/3) = 2p$, em $p \approx 0.6$: a essa altura a
diferença bruta já perdeu a maior parte de sua informação.
\end{solution}

\begin{solution}{exo:b3:molecular-evolution:7}
$p_{N} = 7/700 = 0.010$, $p_{S} = 10/200 = 0.050$, $\omega = 0.2$:
\omterm{def:b3:molecular-evolution:dnds}{seleção purificadora} removendo quatro quintos das mudanças de aminoácido.
Com $35$ e $10$: $p_{N} = 0.05$, $\omega = 1$: nenhuma restrição
detectável --- um \omterm{def:b3:molecular-evolution:duplication}{pseudogene}, ou um gene em que a \omterm{def:b3:molecular-evolution:dnds}{seleção positiva} em
alguns sítios equilibra a purificadora em outros.
\end{solution}

\begin{solution}{exo:b3:molecular-evolution:8}
$k = d/2T = 0.012/(1.3\times 10^{7}) = 9.2\times 10^{-10}$ por sítio por
ano; $\mu = 25k = 2.3\times 10^{-8}$ por sítio por geração; $6\times
10^{9}\times 2.3\times 10^{-8} \approx 140$ mutações novas por genoma
diploide por geração. O sequenciamento direto de pais e filhos encontra
cerca de setenta: a estimativa a partir da divergência é inflada pelo
polimorfismo já presente na população ancestral, que acrescenta algumas
centenas de milhares de gerações à separação aparente.
\end{solution}

\begin{solution}{exo:b3:molecular-evolution:9}
$\mathrm{e}^{-T/2N} = \mathrm{e}^{-1} = 0.37$; fração discordante
$\tfrac{2}{3}
\times 0.37 = 0.25$. Para \qty{10}{\%}:
$\mathrm{e}^{-T/2N} = 0.15$, $T =
2N\ln(1/0.15) = \num{190000}$ gerações.
\end{solution}

\begin{solution}{exo:b3:molecular-evolution:10}
Quando $s \to 0$: $1 - \mathrm{e}^{-2s} \approx 2s$ e $1 - \mathrm{e}^{-4Ns}
\approx 4Ns$, razão $1/2N$. Para $4Ns \gg 1$ o
denominador é $1$ e o numerador, $\approx 2s$. Para $s < 0$ com $4N|s| \gg 1$: numerador $1 -
\mathrm{e}^{2|s|} \approx -2|s|$, denominador $1 - \mathrm{e}^{4N|s|} \approx
-\mathrm{e}^{4N|s|}$, de modo que $P \approx 2|s|\,\mathrm{e}^{-4N|s|}$. Cem vezes abaixo da neutra em $N = 10^{4}$:
$2|s|\,\mathrm{e}^{-4\times 10^{4}|s|} =
5\times 10^{-7}$; $|s| = 1.6\times 10^{-4}$ dá $3.2\times 10^{-4}\times
\mathrm{e}^{-6.4} = 5.3\times 10^{-7}$. Logo $|s| \approx 1.6\times 10^{-4}$, $4N|s| \approx 6.5$: uma desvantagem de um centésimo de por cento basta, em dez mil
indivíduos, para tornar a fixação cem vezes mais rara que o acaso.
\end{solution}

\begin{solution}{exo:b3:molecular-evolution:11}
Com o ponto de separação a uma distância $a$ de $O$ pelos dois caminhos,
$a + m =
0.30$, $a + h = 0.24$, $m + h = 0.20$: $m - h = 0.06$, de modo
que $m = 0.13$, $h
= 0.07$, razão $1.9$. Um relógio estrito por geração
daria a razão das contagens de gerações, $25/0.5 = 50$. O $1.9$ observado
significa que o camundongo acumulou apenas o dobro da mudança humana por
ano, apesar de cinquenta vezes as gerações: a taxa de mutação por geração
precisa ser umas $25$ vezes maior nos humanos --- mais divisões celulares
germinativas por geração ---, de modo que, por ano, as duas linhagens
difiram por um fator de apenas dois.
\end{solution}

\begin{solution}{exo:b3:molecular-evolution:12}
As substituições que criam paradas chegam a $1000\times 2\times 10^{-9}\times
0.04 = 8\times 10^{-8}$ por ano: a primeira é esperada
depois de cerca de \qty{12}{Myr} (as mudanças de fase mais ou menos
reduzem isso à metade). Nessa janela, uma mutação benéfica, uma em
$10^{3}$ a $10^{5}$ de todas as mutações, tem muito menos chance de
chegar que um códon de parada e, mesmo quando chega, só se fixa com
probabilidade $2s$. Assim, a cópia em geral já está morta antes que a
seleção lhe encontre uso --- a meia-vida observada das duplicatas em
vertebrados é de alguns milhões de anos.
\end{solution}

\begin{solution}{pb:b3:molecular-evolution:1}
\textbf{1.} $k = 0.012/(2\times 6.5\times 10^{6}) = 9.2\times 10^{-10}$ por
sítio por ano; $\mu = 25k = 2.3\times 10^{-8}$ por sítio por geração.
\textbf{2.} $6\times 10^{9}\times 2.3\times 10^{-8} \approx 140$ mutações
novas por genoma diploide; \qty{1.5}{\%} delas, cerca de $2$, em
sequência codificadora.
\textbf{3.} Por genoma haploide, $3\times 10^{9}\times 2.3\times 10^{-8}
\approx 70$ substituições fixadas por geração, contra o $1/300$ de
Haldane: vinte mil vezes mais do que a seleção poderia impulsionar, de
modo que a esmagadora maioria é neutra.
\textbf{4.} $2N\mu = 10^{5}\times 2.3\times 10^{-8} = 2.3\times 10^{-3}$
mutações novas por sítio por geração na população; $1/2N =
10^{-5}$ delas
se fixarão.
\textbf{5.} $4N = \num{200000}$ gerações, \qty{5}{Myr} --- quase tanto
quanto o tempo decorrido desde a separação.
\textbf{6.} A divergência conta as mutações que surgiram ao longo das duas
linhagens separadas desde seu ancestral comum; cada uma se fixa dentro de
sua própria linhagem, e a taxa de longo prazo é $\mu$ qualquer que seja o
tempo que cada uma leve, já que as mutações surgem de modo constante e a
demora só atrasa o fluxo, sem alterá-lo. (A única correção é para o
polimorfismo presente na separação, que torna o ancestral comum de dois
alelos mais velho que a separação das espécies.)
\textbf{7.} Neutra, $10^{-5}$; $s = 0.001$: $0.002$; $s = 0.01$: $0.0198$;
$s = 0.1$: $1 - \mathrm{e}^{-0.2} = 0.18$.
\textbf{8.} $s = 1/4N = 5\times 10^{-6}$: abaixo disso a deriva engole a
seleção e a mutação se comporta como neutra; acima disso a seleção
governa as chances.
\textbf{9.} $s = -0.001$: $P = 0.002\,\mathrm{e}^{-200}$, zero para todos
os efeitos --- $10^{-85}$ da neutra. $s = -10^{-5}$, $4N|s| = 2$: $P =
2\times 10^{-5}/(\mathrm{e}^{2} - 1) = 3.1\times 10^{-6}$, \qty{31}{\%} da
neutra: as mutações levemente deletérias de fato se fixam.
\textbf{10.} Substituições benéficas por sítio por geração: $2N\times
10^{-5}\mu\times 2s = 10^{5}\times 10^{-5}\times 0.02\,\mu = 0.02\mu$;
neutras: $\approx\mu$. Fração adaptativa $0.02/1.02 \approx \qty{2}{\%}$.
Embora cada mutação benéfica tenha duas mil vezes mais chance de se fixar
que uma neutra, elas são tão raras que apenas uma substituição em
cinquenta é adaptativa: a seleção age e mal aparece.
\textbf{11.} $p_{N} = 46/920 = 0.050$, $p_{S} = 84/280 = 0.30$.
Corrigidos: $d_{N} = -\tfrac{3}{4}\ln(1 - 0.0667) = 0.052$, $d_{S} = -\tfrac{3}{4}
\ln(1 - 0.40) = 0.38$. $\omega = 0.14$.
\textbf{12.} \omterm{def:b3:molecular-evolution:dnds}{Seleção purificadora} removendo cerca de \qty{86}{\%} das
mudanças de aminoácido --- um gene típico. $k = d_{S}/2T = 0.38/(1.8\times
10^{8}) = 2.1\times 10^{-9}$ por sítio por ano: o dobro do
valor humano--chimpanzé, como se espera de uma comparação que inclui a
linhagem de roedores, de relógio rápido; a mesma ordem de grandeza.
\textbf{13.} $d = -\ln(1 - 0.55) = 0.80$ por sítio; $T = 0.80/(2\times
10^{-9}) = \qty{400}{Myr}$.
\textbf{14.} Duas cadeias diferentes formam o tetrâmero
$\alpha_{2}\beta_{2}$, cuja ligação cooperativa ao oxigênio, \omterm{prop:b3:structural-biology:mechanism}{efeito Bohr}
e controle alostérico pelo 2,3-bisfosfoglicerato dependem da interface
entre subunidades diferentes; e o agrupamento $\beta$ pôde então se
diversificar em cadeias embrionárias, fetais e adultas com afinidades
diferentes.
\textbf{15.} $2kT = 2\times 8\times 10^{-9}\times 4\times 10^{7} = 0.64$
substituição por sítio; com a saturação entre vinte aminoácidos a
diferença bruta é de cerca de $0.95(1 - \mathrm{e}^{-0.64/0.95}) \approx 0.47$. Em \qty{500}{Myr}, $d = 8$: todo sítio já mudou muitas vezes e a
diferença bruta fica em seu teto de \qty{95}{\%}, sem carregar informação
alguma sobre o tempo.
\textbf{16.} $10^{-11}\times 10^{9}\times 100\times 2 = 2$ substituições
entre duas linhagens em 100 sítios ao longo de um bilhão de anos. Para
uma separação de \qty{10}{Myr} a expectativa é $0.02$: quase sempre
nenhuma diferença, nenhuma resolução.
\textbf{17.} $1200\times 2\times 10^{-9}\times 0.04 = 9.6\times 10^{-8}$
por ano: a primeira parada é esperada depois de cerca de \qty{10}{Myr}.
\textbf{18.} Duplicar o genoma inteiro duplica todo gene ao mesmo tempo,
de modo que a estequiometria das proteínas que interagem é preservada e
nenhum desequilíbrio de dosagem seleciona contra as cópias; uma duplicata
isolada em tandem altera a dose de um gene e muitas vezes é deletéria,
ou é imediatamente redundante e livre para decair.
\textbf{19.} $\mathrm{e}^{-\num{80000}/\num{100000}} = \mathrm{e}^{-0.8} = 0.45$.
\textbf{20.} $\tfrac{2}{3}\times 0.45 = 0.30$: consistente com os
\qty{30}{\%} observados.
\textbf{21.} Metade cada, com \qty{15}{\%} de todas as árvores agrupando
humano com gorila e \qty{15}{\%} chimpanzé com gorila. Um claro excesso
de um deles indicaria fluxo gênico entre essas duas espécies após sua
separação, o que a deriva não consegue produzir.
\textbf{22.} $N = \num{10000}$: $\mathrm{e}^{-4} = 0.018$, discordância de
\qty{1.2}{\%}. Os \qty{30}{\%} observados exigem, portanto, uma população
ancestral de uns cinquenta mil --- várias vezes o tamanho efetivo dos
humanos de hoje.
\textbf{23.} A concatenação supõe que todo gene tem a \omterm{prop:b3:molecular-evolution:ils}{árvore de espécies};
com \omterm{prop:b3:molecular-evolution:ils}{separação incompleta de linhagens} os genes têm muitas árvores e, para
ramos internos curtos, a \omterm{prop:b3:molecular-evolution:ils}{árvore gênica} mais comum pode diferir da \omterm{prop:b3:molecular-evolution:ils}{árvore
de espécies}, de modo que mais dados tornam a resposta errada mais
confiante. Um método do coalescente calcula, para cada \omterm{prop:b3:molecular-evolution:ils}{árvore de espécies}
candidata, a distribuição esperada de \omterm{prop:b3:molecular-evolution:ils}{árvores gênicas}, e escolhe a \omterm{prop:b3:molecular-evolution:ils}{árvore
de espécies} que melhor explica as frequências observadas.
\textbf{24.} A \omterm{prop:b3:molecular-evolution:ils}{separação incompleta de linhagens} aconteceu na população
ancestral comum de todos os humanos modernos, de modo que tornaria toda
população humana igualmente aparentada aos neandertais. Um excesso nos
não africanos significa que os alelos neandertais entraram depois que os
ancestrais dos não africanos deixaram a África: mistura, uns cinquenta mil
anos atrás.
\textbf{25.} $\mu \approx 2.3\times 10^{-8}$ por sítio por geração;
$\omega \approx 0.14$; cerca de \qty{30}{\%} das \omterm{prop:b3:molecular-evolution:ils}{árvores gênicas}
discordantes da \omterm{prop:b3:molecular-evolution:ils}{árvore de espécies}.
\end{solution}
